% ЕМТ 2022, X клас — проверен LaTeX препис на липсващата задача 10.1 от официалните файлове в Klasirane.com.
% Условия: https://klasirane.com/api/competitions/EMT/All/2022/10%20%D0%BA%D0%BB/probs
% SHA-256 (условия): aeae6e35515428783ccfb282c49aa8f467340c3963b531ca30eeff3b278b900f
% Решения: https://klasirane.com/api/competitions/EMT/All/2022/10%20%D0%BA%D0%BB/sol
% SHA-256 (решения): b4593a61b3cbcdbcadda13d44594ad318f1fb7e860d93fcf594b7ff9bf435cb0
% Mathpix Snip (условия): b38a8027-f697-4e9a-bd05-24ca5fc4f01e
% Mathpix Snip (решения): d2729e35-ea97-4097-881c-5478c74c4d68
% Формулите, формулировките и точковата оценка са сверени визуално с официалните PDF страници.
% В официалните файлове няма фигура към задача 10.1.

Задача 10.1. Да се реши ирационалното уравнение:
$$
3 \sqrt{3 x-1}=x^{2}+1 .
$$

Решение. Отговор: $x_{1,2}=\frac{3 \pm \sqrt{5}}{2}$.

Даденото уравнение има смисъл при $x \geq \frac{1}{3}$. Нека въведем ново неизвестно $t=\sqrt{3 x-1} \geq 0$. Тогава ирационалното уравнение е еквивалентно на следната система:
$$
\left\{
\begin{aligned}
t &\geq 0, \\
t^{2} &=3 x-1, \\
3 t &=x^{2}+1
\end{aligned}
\right.
\quad \text{или} \quad
\left\{
\begin{aligned}
t &\geq 0, \\
3 x &=t^{2}+1, \\
3 t &=x^{2}+1
\end{aligned}
\right. \tag{*}
$$
(За решенията на тази система ще бъде изпълнено и $x \geq \frac{1}{3}$.)

Сега след почленно изваждане на двете уравнения на системата (*) последователно получаваме $3(x-t)=t^{2}-x^{2}$, $(x-t)(3+x+t)=0$, т.е. $x=t$ или $x+t+3=0$.

Следователно системата (*) е еквивалентна на обединението на двете системи
$$
\left\{
\begin{aligned}
t &\geq 0, \\
x &=t, \\
3 t &=x^{2}+1
\end{aligned}
\right. \tag{**}
\qquad \text{и} \qquad
\left\{
\begin{aligned}
t &\geq 0, \\
x+t+3 &=0, \\
3 t &=x^{2}+1
\end{aligned}
\right. \tag{***}
$$

След заместване на $t=x$ във второто уравнение на (**) достигаме до $x^{2}-3 x+1=0$, чиито решения са $x_{1}=\frac{3+\sqrt{5}}{2} \geq 0$ или $x_{2}=\frac{3-\sqrt{5}}{2} \geq 0$. Така получаваме, че решенията на (**) са
$$
\left\{
\begin{aligned}
x_{1}&=\frac{3+\sqrt{5}}{2}, \\
t_{1}&=\frac{3+\sqrt{5}}{2}
\end{aligned}
\right.
\quad \text{или} \quad
\left\{
\begin{aligned}
x_{2}&=\frac{3-\sqrt{5}}{2}, \\
t_{2}&=\frac{3-\sqrt{5}}{2}.
\end{aligned}
\right.
$$

Заместваме $t=-3-x$ във второто уравнение на (***) и свеждаме до $x^{2}+3 x+10=0$, което няма реални корени, т.е. системата (***) също няма реални корени.

Сега окончателно получаваме, че решенията на ирационалното уравнение $3 \sqrt{3 x-1}=x^{2}+1$ (системата (*)) са $x_{1}=\frac{3+\sqrt{5}}{2}$ или $x_{2}=\frac{3-\sqrt{5}}{2}$.

Оценяване. (6 точки) 1 точка за дефиниционното множество. 1 точка за полагането. 1 точка за $x=t$ и $x+t+3=0$ и достигане до (**) и (***). 2 точки за намиране на решенията на (**) и (***). 1 точка за окончателен отговор.

Алтернативно: (6 точки) За трансформиране на уравнението в полином от четвърта степен и разлагането му на два квадратни полинома — 4 точки. За решаване на всеки от квадратните полиноми — по 1 точка.
